\par\noindent\textbf{(1)} Zeeman分裂使两自旋能量相对偏移 \(\mp\mu_BB\)，因而磁化 \(M=\mu_B\int d\varepsilon\,g_s(\varepsilon)[f_F(\varepsilon-\mu_BB)-f_F(\varepsilon+\mu_BB)]\)。对小场展开占有差为 \(-2\mu_BB f_F'(\varepsilon)+O(B^3)\)。
\par\noindent\textbf{(2)} 于是线性响应 \(\chi_B=\partial M/\partial B=2\mu_B^2\int_0^\infty g_s(\varepsilon)[-f_F'(\varepsilon)]d\varepsilon\)。零温时 \(-f_F'(\varepsilon)\to\delta(\varepsilon-E_F)\)，故 \(\chi_B(0)=2\mu_B^2g_s(E_F)\)，是费米面附近态密度而非全部电子数直接贡献。
\par\noindent\textbf{(3)} 三维抛物色散总态密度 \(g_{\rm tot}(E_F)=3N/(2E_F)\)，单自旋为 \(g_s=3N/(4E_F)\)。因此样品总磁化对 \(B\) 的响应为 \(\chi_B=3N\mu_B^2/(2E_F)\)；若取体积磁化 \(M/V\)，应再除以 \(V\)，对外场 \(H\) 的SI磁化率还需相应乘 \(\mu_0\)。
\par\noindent\textbf{(4)} 有限低温用Sommerfeld展开，磁响应主项为费米面态密度，修正 \(O[(k_BT/E_F)^2]\)。经典非简并高温极限则 \(\chi_B\propto N\mu_B^2/(k_BT)\) 呈Curie律，因Fermi–Dirac占有不再阶跃；不可把零温Pauli常数外推到高温。
